%%% Methodology --- 
%% 
%% Filename: Method.tex
%% Author: Jacob Ringfjord & Max Wilén
%%% Code:

\chapter{Methodology}
\label{cha:methodology}
\begin{comment}
OLD:
In Section~\ref{sec:contributions}, the overarching methodological approach was outlined. This section builds on that foundation, providing a more detailed elaboration, with each step being covered individually. The chosen methodology is supported by insights from prior research in the field of computer science and software engineering, along with a discussion on how the result can prove valuable and applicable in terms of the project's context.
\end{comment}

The research for this thesis began by exploring relevant tools for binary comparison. Numerous options exist as outlined in Section~\ref{sec:background:binary_disassembler_decompilers} and Section~\ref{sec:background:binary_diffing}. Experimentation and comparison led to selecting a tool that supports multiple architectures and large-scale programs, and one that meets this thesis's current needs and the client's future needs when application to other architectures and functionality extendability is requested. This reasoning led us to choose Ghidriff as the solution. It leverages the power of the well-established Ghidra engine and presents detailed binary diffing output for use during change categorization.


%\section{Experimental development}
%This section of the methodology %concerns the experimental %development of the tool this %thesis presents, and the %methodological steps that made %the tool. 

\label{sec:method:overview}
The pipeline developed in this thesis can be split into three separate modules. The first covers the comparison of the two given executables. This module is named \textit{Change detection} and uses Ghidra to produce the binary diffing outputs needed for the second module, named \textit{Categorization}, which categorizes the detected changes. The third module \textit{Evaluation} evaluates these categorized changes alongside their relevance. The output from these three modules is later presented in an HTML report that ranks the changes found based on the evaluated relevance and also presents the specific logical functionality that has changed between the two executables. This overview is presented in Figure~\ref{fig:pipeline_overview2} and resembles the flow of the developed tool in a high-level manner.

\begin{figure}[ht]
    \centering
    \begin{dot2tex}[dot, options=-tmath --autosize, scale=0.8]
    digraph G {
        rankdir=LR;
        node [shape=box, style=filled, fillcolor=lightgray];

        "\text{Old APK}" -> "\text{Ghidriff}";   
        "\text{New APK}" -> "\text{Ghidriff}";

        // Subgraph for Input & Ghidriff
        subgraph cluster_input {
            label = "\text{Change detection}";
            style = rounded;
            "\text{Ghidriff}";
        }

        // Subgraph for Categorization
        subgraph cluster_categorization {
            label = "\text{Categorization}";
            style = rounded;
            "\text{Ghidriff}" -> "\text{Syntax matcher}";
        }

        // Subgraph for Ranking
        subgraph cluster_ranking {
            label = "\text{Evaluation}";
            style = rounded;
            "\text{Syntax matcher}" -> "\text{Weighted words}";
        }

        "\text{Weighted words}" -> "\text{Result}";
    }
    \end{dot2tex}
    \caption{An overview of the different modules in the developed tool}
    \label{fig:pipeline_overview2}
\end{figure}



\section{Dependencies and Technical Prerequisites}
The pipeline is written in Python and fetches JSON data produced by Ghidriff to perform the classification. During classification, data is stored inside the running pipeline program as lists of functions, symbols, etc. in a typical Python manner. The evaluation and relevance prioritization are performed on the same data. After the evaluation is complete, an HTML report is produced by the pipeline, which visualizes the findings. As mentioned, this report is visualized in the screenshot in Figure~\ref{fig:pipeline_frontpage}. See Figure~\ref{fig:project_structure} for project structure.


\begin{figure}[H]
    \centering
    \resizebox{0.5\textwidth}{!}{%
        \begin{forest}
            for tree={font=\sffamily, 
            grow'=0,
            folder indent=.9em, 
            folder icons,
            edge=densely dotted}
            [\
                [categorization
                    [categorization.py, is file]
                    [metadata\_registration.py, is file]
                    [matchers
                        [\_\_init\_\_.py, is file]
                        [<CATEGORY\_NAME>.yml, is file]
                    ]
                    [helpers.py, is file]
                    [static.py, is file]
                ]
                [change\_detection
                    [ghidriff\_controller.py, is file]
                ]
                [evaluation
                    [eval.py, is file]
                ]
                [misc
                    [FuncCollection.py, is file]
                    [FuncNode.py, is file]
                    [FuncStatus.py, is file]
                    [GhidriffLog.py, is file]
                    [report\_generation.py, is file]
                    [static.py, is file]
                ]
                [signal-binaries]
                [ghidriff-project]
                [ghidriffs]
                [output]
                [result
                    [X.XX.X-Y.YY.Y
                        [analysis-report.html, is file]
                        [categories\_critical.html, is file]
                        [categories\_noncritical.html, is file]
                        [function\_data.json, is file]
                        [overview.html, is file]
                    ]
                ]
                [pipeline.py, is file]
                [run.bat, is file]
                [run.sh, is file]
                [requirements.txt, is file]
                [README.md, is file]
            ]
        \end{forest}
    }
    \caption{Project structure}
    \label{fig:project_structure}
\end{figure}



\section{Change Detection}

The tool Ghidriff, described in Section~\ref{sec:background:ghidriff}, is chosen to detect changes between the software versions. Compared to the pipeline mentioned in Section~\ref{sec:related_work:kandidat_pipeline}, the current implementation of Ghidriff can neither categorize the changes nor rank them. However, the result of binary analysis and binary diffing is much more informative and serves as a better platform of data to expand upon than the pipeline described in Section~\ref{sec:related_work:kandidat_pipeline}. Because of this and the sparse execution time for this thesis project, potential expandability of a tool without the need to change the tool itself is seen as an advantage when deciding which tool to build upon. Because of this, Ghidriff has been chosen over the two.

Other tools for detecting changes have been described in Sections~\ref{sec:background:binary_disassembler_decompilers}~and~\ref{sec:background:binary_diffing}. E.g., recall that combining IDA Pro and BinDiff can serve as a useful alternative to Ghidriff as the combination leverages two independently strong engines for binary analysis and binary diffing. However, integrating these two into a pipeline would require additional time to ensure they work together effectively, rather than simply using Ghidriff's output data. Ultimately, the decision comes down to leveraging an already robust and extendable platform rather than building a new one. 
% More about this can be found in Section~\ref{sec:discussion-method}.

In Figure~\ref{fig:ghidriff_diagram}, an in-depth visualization of how Ghidriff is constructed is presented, and examples of what the JSON output looks like are found in Figures~\ref{fig:ghidriff_json_output_A},~\ref{fig:ghidriff_json_output_B},~and~\ref{fig:ghidriff_json_output_C}.

%...The categorization used in the work presented by %\textcite{nfc_kandidatrapport_jesper}, mentioned in Section %\ref{sec:related_work:kandidat_pipeline} is in early development. Because of %this, in tandem with the lack of data, their categorization computes on, %Ghidriff is the chosen option of the two and the work that this thesis %builds towards.

\tikzstyle{startstop} = [rectangle, rounded corners, minimum width=5cm, minimum height=1.5cm, text-centered, draw=black, fill=color1, drop shadow]
\tikzstyle{process} = [rectangle, rounded corners, minimum width=5.5cm, minimum height=2cm, text-centered, draw=black, fill=color0, drop shadow]
\tikzstyle{version_comp} = [rectangle, rounded corners, minimum width=5.5cm, minimum height=2cm, text-centered, draw=black, fill=color3, drop shadow]
\tikzstyle{diffing} = [rectangle, rounded corners, minimum width=5.5cm, minimum height=2cm, text-centered, draw=black, fill=color2, drop shadow]
\tikzstyle{ours} = [rectangle, rounded corners, minimum width=5.5cm, minimum height=2cm, text-centered, draw=black, fill=white!20, drop shadow]
\tikzstyle{arrow} = [thick,->,>=stealth]
\tikzstyle{dashed_arrow} = [dashed,->,>=stealth]
\begin{figure}
    \centering
    \begin{tikzpicture}[
        node distance=1cm and 1cm, 
        every node/.style={font=\footnotesize}
    ]   

        % Nodes with bullet lists
        \node (start) [startstop] {%
            \textbf{Input File}
            \begin{tabular}{l}
                $\bullet$ Prepare two binaries ("new" and "old") \\
            \end{tabular}
        };

        \node (init) [process, below=of start] {%
            \textbf{Initialize Ghidra}
            \begin{tabular}{l}
                $\bullet$ Launch headless version of Ghidra (using ProgramAPI and FlatProgramAPI) \\ 
                $\bullet$ Load binaries into Ghidra \& create Ghidra project structure
            \end{tabular}
        };


        % SECTION DIVIDER 1
        \draw[dashed, thick] ($(init.south) + (-8,-0.5)$) -- ($(init.south)+ (6,-0.5)$);
        \node at ($(init.north) + (-6.7,0.5)$) {\textbf{Setup Ghidra engine}};


        \node (control) [version_comp, below=of init] {%
            \textbf{Analyze binaries}
            \begin{tabular}{l}
                $\bullet$ Ghidra performs separate analysis on binaries provided \\ 
                $\bullet$ Disassembly and decompilation of each binary \\
                $\bullet$ Extraction of metadata, symbols, functions, libraries, etc.
            \end{tabular}
        };

        \node (parse) [version_comp, below=of control] {%
            \textbf{Store analysis data}
            \begin{tabular}{l}
                $\bullet$ Store findings in temporary Ghidra \\project structure for later diffing
            \end{tabular}
        };


        % SECTION DIVIDER 2
        \draw[dashed, thick] ($(parse.south) + (-8,-0.5)$) -- ($(parse.south)+ (6,-0.5)$);
        \node at ($(parse.north) + (-6.7,0.5)$) {\textbf{Version comparison}};


        \node (match) [diffing, below=of parse] {%
            \textbf{Identify changes}
            \begin{tabular}{l}
                $\bullet$ Compare function signatures \\ 
                $\bullet$ Find structural matches \\ 
                $\bullet$ Match similar functions using various correlators \\ 
            \end{tabular}
        };

        \node (diff) [diffing, below=of match] {%
            \textbf{Parse changes}
            \begin{tabular}{l}
                $\bullet$ Detect added \& deleted functions \\ 
                $\bullet$ Detect added \& deleted symbols \\ 
                $\bullet$ Track differences in modified functions
            \end{tabular}
        };

        % SECTION DIVIDER 3
        \draw[dashed, thick] ($(diff.south) + (-8,-0.5)$) -- ($(diff.south)+ (6,-0.5)$);
        \node at ($(diff.north) + (-6.7,0.5)$) {\textbf{Binary diffing}};

        \node (output) [startstop, below=of diff] {%
            \textbf{Output File(s)}
            \begin{tabular}{l}
                $\bullet$ Generate JSON output \\ 
                $\bullet$ Generate Markdown reports \\ 
                $\bullet$ Changes in modified decompiled code
            \end{tabular}
        };

        %\node (extension) [ours, below=of output] {%
        %    \textit{Our pipeline extension}
        %};



        % Arrows
        \draw [arrow] (start) -- (init);
        \draw [arrow] (init) -- (control);
        \draw [arrow] (control) -- (parse);
        \draw [arrow] (parse) -- (match);
        \draw [arrow] (match) -- (diff);
        \draw [arrow] (diff) -- (output);
        %\draw [dashed_arrow] (output) -- (extension);

    \end{tikzpicture}
    \caption{Inner workings and flow of Ghidriff}
    \label{fig:ghidriff_diagram}
\end{figure}


\section{Categorization}
\label{sec:method:categorization}
\begin{comment}
As mentioned earlier the main topic of this thesis is to categorize the changes based on the result from the binary analysis and binary diffing. To achieve this machine learning will be used and integrated with the Ghidriff tool. 

What to categorize?
how important/severe is the change?
What type of change --> from key changes?
"reason" for change, security patch, bug fix, optimization, new feature etc.
\end{comment}

When performing the change detection on released software applications, the output will most likely be of such size that manually looking at the result will prove inefficient. As explained in Section~\ref{sec:background:compilation_effects_on_binaries},~and~Section~\ref{sec:background:comparing_binaries}, the wide range of configurations and compilation options can result in numerous irrelevant or superficial differences being detected. Because of this, it is essential to implement automated mechanisms for categorizing changes based on their significance, ensuring that only meaningful or potentially impactful modifications are flagged for further analysis. Because this thesis is focused on providing a valuable tool for use in a forensic law enforcement process, as well as focusing on Android applications, the categorization is based on Android-specific operations that alter the logical, functional, and structural flow of the application.

\subsection{Changes to Categorize}
\label{sec:method:changes_to_categorize}
For this thesis, the client provided an initial list of interesting changes to categorize. After further discussions, the list was polished, and the critical changes were defined. These changes are considered to be of extra importance and relevance in a forensic context and can be seen in Table~\ref{tab:critical_changes}. From here onwards, these types of changes are referred to as \textit{critical categories}.

\begin{table}[H]
    \centering
    \begin{tabular}{p{3cm} p{10cm}}
        \toprule
        \textbf{Type of change} & \textbf{Details} \\
        \midrule
        
        \codebox{Data storage} &     
        Changes in data storage structure, new reads and writes, etc. \\\\
        \midrule
        \codebox{File write} & 
        Changes resulting in file write additions or deletions of files.\\\\
        \midrule
        \codebox{Logging} & 
        Changes in what, where, when, and how data is logged. \\\\
        \midrule
        \codebox{Authentication} & 
        Changes in login and logout processes, session handling, token validation, and other authentication mechanisms. \\\\
        \midrule
        \codebox{Error handling} & 
        Adjustments to how errors and exceptions are caught, handled, or logged, including retries, rollbacks, and fallback mechanisms. \\\\
        \midrule
        \codebox{Cryptographic} & 
        Changes to encryption algorithms, key management, hashing methods, etc.\\\\
        \bottomrule
    \end{tabular}
    \caption{Critical software changes to identify for forensic analysis}
    \label{tab:critical_changes}
\end{table}



To be certain in finding these critical changes, the categorization also looks at changes that are considered non-critical. This minimizes the risk of overlooking modifications that might not directly match predefined critical patterns but could still impact core functionalities. See Table~\ref{tab:noncritical_changes} for typical non-critical operations for Android applications, which can be considered critical when analyzed structurally. From here onwards, these types of changes are referred to as \textit{non-critical categories}.

\begin{table}[H]
    \centering
    \begin{tabular}{p{3cm} p{10cm}}
        \toprule
        \textbf{Type of change} & \textbf{Details} \\
        \midrule
        
        \codebox{UI update} &     
        Changes related to user interface components, such as modifications to views, layouts, colors, visibility toggles, and animations. \\\\
        \midrule
        \codebox{Network call} & 
        Alterations in network requests, including changes in API calls, request/response handling, and new or removed endpoints. \\\\
        \midrule
        \codebox{Library function} & 
        Changes to functions originated from external or built-in library packages, rather than self-developed ones. \\\\
        \midrule
        \codebox{Text manipulation} & 
        Modifications to text strings, including changes to messages, labels, prompts, or other static textual content. \\\\
        \bottomrule
    \end{tabular}
    \caption{Non-critical software changes for forensic analysis}
    \label{tab:noncritical_changes}
\end{table}


The categorization aims to streamline the process of identifying interesting changes between software versions and reducing reliance on ad-hoc methods. Categorizing changes into multiple categories is, therefore, not considered an issue since it still fits the purpose. Changes affecting more than one thing are also not rare occurrences, and categorizing these into more than one category is logical.

%\textbf{TODO:} Add 1-2 sentences stating how the categories were selected (e.g., discussions with NFC experts, review of forensic case studies).


\subsection{Syntax Matching}
By analyzing code changes at the syntactic level, it is possible to leverage pattern-matching techniques to detect specific operations such as the ones listed in Table~\ref{tab:critical_changes}. This can be done by using regular expressions to match commonly used symbols and function names to understand what operations are most likely occurring. Adding to this, analyzing structural patterns and how these symbols and functions are interlinked is also important.

The core process of categorizing the changes begins with parsing the binary diffing output from Ghidriff. Here, the function, library, and superclass names are shown, as well as information about parent and child functions in the call graph. Adding to this, changes in the decompiled bytecode, which Ghidra decompiles into representative pseudo code, can also be parsed in this step. Once the diff is processed, regular expressions are applied to identify key syntactic patterns corresponding to predefined categories. 




\subsubsection{Android Syntax Matching}
The syntax matching of Android-specific operations used in this thesis is based on the assumption that industry-standard function names and symbol patterns serve as effective indicators of program behavior. The set of keywords is selected due to their widespread usage in software development related to Android frameworks, APIs, and libraries, where they appear consistently across applications. It is assumed that developers generally use industry-standard functions over custom implementations for reasons related to security, reliability, and ease of use. It is also believed that developers rarely create entirely new function names for common tasks and rely on these established naming conventions, making this approach feasible.




\subsubsection{Categorization Using Regex}
Regular expression grammars have been defined for each category to map function names, symbols, and code changes accordingly. During syntax analysis, function names are compared against a predefined set of keywords. In addition, critical categories are also computed using an analysis of the code diff. The reason for only looking at the code diff when computing the critical syntax analysis is that the critical changes are harder to spot when only matching function names. This is further discussed in Section~\ref{sec:discussion-method}

To illustrate the syntax analysis, consider a function that has been identified as changed between the two compared binaries. Then, the function names and code changes will undergo a syntax matching analysis where the strings are mapped to the predefined list of keywords for each of the available categories. All matched keywords are stored for later use when generating the HTML report, and the functions are assigned to the matched category. The keywords for the regex matching used in this thesis are found in the Tables~\ref{tab:critical_regex_keywords}~and~\ref{tab:noncritical_regex_keywords}.\textbf{}

\subsubsection{Two-layered Matching}
\label{sec:two_layered_matching}
The keywords used for critical changes in the syntax matching undergo a two-layered matching process made by us, where the words are divided into two types: \textit{normal} and \textit{ambiguous}. \textit{Normal} words are more reliably linked to a specific category compared to the \textit{ambiguous} ones. For instance, \textit{db} is more likely to yield an accurate categorization than \textit{delete}, as the latter is assumed to be used more often in other contexts than only database management. The sectioning between the two types for critical keywords is presented in  Table~\ref{tab:critical_regex_keywords}.

However, this separation does not change the behavior of the categorization, only the evaluation. For instance, the categorization for a change containing the word \textit {delete} should be considered as a Data storage change, but when comparing it with other changes of the same category, it should be ranked lower than, let's say \textit{db}.

\begin{table}[H]
  \centering
  \begin{tabular}{@{} l l l l l l @{}}
    \toprule
    \codebox{Authentication} & \codebox{Cryptographical} & \codebox{Data Storage} & \codebox{Error Handling} & \codebox{File Write} & \codebox{Logging} \\ \midrule
    auth           & aes               & commit        & assert         & append*     & Log\_d \\
    bearer         & brsa              & database      & catch          & export      & Log\_e \\
    credentials    & bouncycastle      & db            & crash          & file        & Log\_i \\
    jwt            & cipher            & delete*       & error          & persist     & Log\_v \\
    login          & Cipher            & insert*       & exception      & save*       & Log\_w \\
    logout         & crypto            & query*        & fail           & write*      & logger \\
    MasterKey      & CryptoHelper      & session*      & handle*        &             & log \\
    mfa            & CryptoUtils       & sql           & recover*       &             & print* \\
    oauth          & decrypt           & store*        & retry*         &             & trace \\
    password       & digest            & transaction   & rollback       &             & verbose \\
    refresh*       & Encrypted         &               & throw          &             & warning \\
    session*       & Encryptor         &               &                &             & warn \\
    signin         & hash*             &               &                &             & wtf \\
    signout        & hmac              &               &                &             & debug* \\
    signup         & java.security     &               &                &             & \\
    sso            & javax.crypto      &               &                &             & \\
    token          & keystore          &               &                &             & \\
    validate       & keygen            &               &                &             & \\
    verify         & KeyStore          &               &                &             & \\
                  & MessageDigest     &               &                &             & \\
                  & md5               &               &                &             & \\
                  & SecretKey         &               &                &             & \\
                  & SecureStore       &               &                &             & \\
                  & sha256            &               &                &             & \\
                  & signature         &               &                &             & \\
                  & spongycastle      &               &                &             & \\ 
                  & verify*           &               &                &             & \\
\bottomrule
  \end{tabular}
  \caption{Keywords used for critical categories. *Word of type ambiguous}
  \label{tab:critical_regex_keywords}
\end{table}

\begin{table}[H]
  \centering
  \begin{tabular}{@{} l l l @{}}
    \toprule
    \codebox{Network Call} & \codebox{UI Update}       & \codebox{Text Manipulation} \\ \midrule
    api          & animatingtoggle & replace \\
    connect      & bui             & spannable \\
    fetch        & compact         & string \\
    http         & display         & substring \\
    network      & font            & text \\
    post         & fragment        &  \\
    put          & hide            &  \\
    receive      & inflate         &  \\
    request      & padding         &  \\
    send         & refresh         &  \\
    socket       & render          &  \\
    sync         & scaffold        &  \\
                 & setvisibility   &  \\
                 & show            &  \\
                 & toggle          &  \\
                 & update          &  \\
                 & view            &  \\
                 & windowsize      &  \\ \bottomrule
  \end{tabular}
  \caption{Keywords used for non-critical categories}
  \label{tab:noncritical_regex_keywords}
\end{table}





\subsubsection{Syntax Classification Errors}
\label{sec:method:false_pos_false_neg}
A \textit{false positive} occurs when a classification incorrectly determines a result to be true when it is false. A \textit{false negative} occurs when the classification fails to detect a true result. Here follows further explanation in regards to this thesis.

\begin{description}
    \item[False positives] When the regex matcher either \textit{(1)} matches a substring in a longer word and incorrectly categorizes it, or \textit{(2)} matches a valid word that appears in a context unrelated to the category to which the word belongs.

    \item[False negatives] When the regex matching does not categorize a change that should be categorized. This is due to missing keywords in the regex matcher, since it only detects changes for the pre-defined keywords.
\end{description} 


Incorrect or misleading variable names can indeed affect the categorization and lead to false positives and false negatives, making these errors challenging to detect. The regex matcher is sound in that its categorization is strictly based on the predefined keywords used in the regex pattern. However, to achieve completeness would require every possible change to be categorized, which is practically infeasible. Therefore, a fallback category for uncategorized changes is used as a compromise.

\section{Evaluation}
%Comparing versions in git logs / Retrieving data changes from
%looking at the git repo
The evaluation of the categorized changes is performed by weighting collected strong and weak words mentioned in Section~\ref{sec:two_layered_matching}, alongside with the amount of assigned categories for that change. Changes with the highest scores are considered to be the most important and therefore ranked the highest. The aim is to ensure that higher-ranked changes are prioritized to effectively highlight the important changes.


\section{Present Results}
The categorization result is parsed in tandem with the analysis data into an HTML report, listing everything needed to make a proper assumption of what changes have occurred in the code base. See Figure~\ref{fig:pipeline_frontpage} for a screenshot of the generated HTML report. This report is explained in detail in Section~\ref{sec:result:pipeline_analysis_report}.

